\documentclass[10pt,a4paper]{amsart}
\usepackage[margin=19mm]{geometry}
\usepackage{amsmath,amssymb,mathtools}
\usepackage[scr=rsfs,cal=euler]{mathalfa}
\usepackage{xfrac}
\usepackage[unicode,hidelinks,bookmarks=false]{hyperref}
\newcommand{\ie}{i.e.\ }
\newcommand{\cf}{cf.\ }
\newcommand{\resp}{respectively\ }
\newcommand{\Subsection}{\subsection}
\providecommand{\nobreakdash}{\nobreak\hbox{-}}
\setlength{\emergencystretch}{2em}
\multlinegap0pt
\usepackage{bm}
\usepackage{bbm}
\usepackage{dsfont}
\usepackage{xspace}
\usepackage{cancel}
\usepackage{mathtools}
\usepackage{amsthm}
\makeatletter
\@ifundefined{thm}{\newtheorem{thm}{Thm}}{}
\@ifundefined{proposition}{\newtheorem{proposition}[thm]{Proposition}}{}
\@ifundefined{theorem}{\newtheorem{theorem}[thm]{Theorem}}{}
\makeatother
\makeatletter
\expandafter\ifx\csname varthm\endcsname\relax
\newenvironment{varthm}[1]{\def\varthmname{#1}\begin{thevarthm}}{\end{thevarthm}\def\varthmname{}}
\fi
\makeatother
\title[On free line arrangements with double, triple and quadruple ]{On free line arrangements with double, triple and quadruple points}
\author{Marek Janasz, Izabela Le\'sniak}
\date{Source: arXiv:2504.05872v2 (CC BY 4.0)}
\begin{document}
\maketitle

\noindent\emph{Source and licence.} On free line arrangements with double, triple and quadruple points. Version arXiv:2504.05872v2 (CC BY 4.0), distributed under the Creative Commons Attribution 4.0 International licence, as recorded in the arXiv licence metadata for this exact version and shown on the abstract page https://arxiv.org/abs/2504.05872v2. Full rights notice: licenses/arxiv-2504.05872-CC-BY-4.0.txt.\par\smallskip

\noindent\textit{Editorial scope.} Counted unit: the Proposition whose source label is \texttt{EstimDLine} in the article (source0.tex lines 158--169, source label \texttt{EstimDLine}), delivered together with the article's own complete original proof of it (source0.tex lines 170--201, one \texttt{proof} environment of 32 lines). No mathematics is written, repaired or re-derived: statement and proof are the article's own text line for line. Required context retained uncounted: dup (lines 127--138). Every definition, display and label of those lines is retained unchanged. Omitted from this package and not redistributed: 1--126, 139--157, 202--571. Nothing inside a retained excerpt is omitted or altered: every retained body is delivered exactly as the article's lines are written, so no omission receipt of any kind accompanies this package. Citations: the retained excerpts carry no citation command at all, so no bibliography entry of the article is transcribed and no reference is redistributed. Reference adaptation: none was needed and none was made; every reference inside the delivered unit resolves inside it, so no unresolved reference or citation is produced. Printed slips kept literal: no printed word, symbol, hypothesis or number of the counted statement or of its proof was altered. Layout and class: the article's own class is \texttt{article} with packages including geometry, hyperref, amsfonts, amssymb, amsthm, amsmath, eucal, tabu, url, pgf, array, tikz-cd, pstricks, pstricks-add; the wrapper is the lane's pinned \texttt{amsart} editorial wrapper at 10pt on A4, which restates the theorem-like environments the retained text needs and adds the article's own macro definitions that the retained text needs (none), with extra wrapper packages (bm, bbm, dsfont, xspace, cancel, mathtools). The delivered numbering is the unit's own. Assets: no retained span contains a figure environment, a graphics-inclusion command, a PDF-page inclusion, a table or a TikZ picture; no image file is copied into this package, the package declares no asset dependency and no image file is redistributed with it. Licence: version arXiv:2504.05872v2 of this article is distributed under the Creative Commons Attribution 4.0 International licence, as recorded in the arXiv licence metadata for this exact version and shown on the abstract page https://arxiv.org/abs/2504.05872v2. A full rights notice is delivered at licenses/arxiv-2504.05872-CC-BY-4.0.txt. Characters: every retained excerpt is pure ASCII, so the delivered engine, XeLaTeX with its default Latin Modern text font, reports no missing character.
\begin{theorem}[du-Plessis -- Wall]
\label{dup}
Let $C = \{f=0\}$ be a reduced plane curve of degree $d$ and let $r = {\rm mdr}(f)$. Let us denote the total Tjurina number of $C$ by $\tau(C)$.
Then the following two cases hold.
\begin{enumerate}
\item[a)] If $r < d/2$, then $\tau(C) \leq \tau_{max}(d,r)= (d-1)(d-r-1)+r^2$ and the equality holds if and only if the curve $C$ is free.
\item[b)] If $d/2 \leq r \leq d-1$, then
$\tau(C) \leq \tau_{max}(d,r)$,
where, in this case, we set
$$\tau_{max}(d,r)=(d-1)(d-r-1)+r^2- \binom{2r-d+2}{2}.$$
\end{enumerate}
\end{theorem}

\begin{proposition}\label{EstimDLine}
Let $\mathcal{L} = \{f = 
0\}\subset\mathbb{P}^2_{\mathbb{C}}$ be a 
free line arrangement of $d$ 
lines with  at 
most quadruple intersection points. 
Then 
\begin{align}\label{EstimateD}
n_2+n_3\leq \bigg\lfloor\frac{3d-3}{2}\bigg\rfloor \:\:\:\text{and}\:\:\:  
\bigg\lceil\frac{d^2-10d+9}{12}\bigg\rceil\leq n_4.     
\end{align}
\end{proposition}

\begin{proof}
Recall that the following naive combinatorial count holds for $\mathcal{L}$:
$$\frac{d(d-1)}{2}=\binom{d}{2}=n_2+3n_3+6n_4.$$
Furthermore, the total Tjurina number of $\mathcal{L}$ is equal to
$$\tau(\mathcal{L}) = n_{2} + 4n_{3} + 9n_{4}.$$
Combining these two facts, we get
$$d(d-1)=\tau(\mathcal{L})+n_2+2n_3+3n_4.$$
Since $\mathcal{L}$ is free, by Theorem \ref{dup} we have
$$r^2-r(d-1)+(d-1)^2=\tau(\mathcal{L}),$$
where $r={\rm mdr}(f)$.
Observe that we can rewrite the above identity as
$$r^2-r(d-1)+(d-1)^2=d(d-1)-n_2-2n_3-3n_4,$$
and hence we get
$$r^2-r(d-1)-d+1+n_2+2n_3+3n_4=0.$$
We compute now the discriminant $\triangle_{r}$ for the above quadratic equation and we get
$$\Delta=d^2-2d+1+4d-4-4n_2-8n_3-12n_4.$$
The freeness of $\mathcal{L}$ implies
$$d^2+2d-3-4n_2-8n_3-12n_4 \geq 0.$$
Using again the naive combinatorial count we obtain
$$d^2+2d-3 \geq 4n_2+8n_3+12n_4 \geq 2(n_2+n_3)+d(d-1),$$
and hence
$$2(n_2+n_3) \leq 3d-3,$$
which gives us first estimate. Since $n_{2}+n_{3}$ is an integer, then we can take the flooring. 

Let us now focus on the lower bound on the number of quadruple intersections. 
Using the naive combinatorial count and our previous bound $n_{2}+n_{3} \leq (3d-3)/2$, we get

$$d^2-d=2n_2+6n_3+12n_4 \leq 6(n_2+n_3)+12n_4 \leq 9d-9+12n_4$$
and hence
$$n_4 \geq \frac{d^2-10d+9}{12}.$$
Since $n_{4} \in \mathbb{Z}_{+}$, we can take the ceiling. Obviously the above bound is non-trivial provided that $d>10$.
\end{proof}

\end{document}
